\section{Thử nghiệm mở rộng với các mô hình và cơ chế học cho dữ liệu mất cân bằng}

Sau các thử nghiệm tái hiện và cải tiến dựa trên XGBoost, nhóm tiếp tục khảo sát khả năng cải thiện bằng cách thay đổi họ thuật toán và cơ chế xử lý mất cân bằng. Mục tiêu không chỉ là tăng Accuracy, mà ưu tiên khả năng phát hiện lớp \textit{stroke} trong khi hạn chế cảnh báo dương tính giả. Các phương pháp được khảo sát gồm CatBoost, LightGBM, Balanced Random Forest, EasyEnsemble và XGBoost với Focal Loss. Đây là ba hướng mở rộng khác nhau: thay đổi họ boosting, xử lý mất cân bằng ở cấp độ thuật toán và sử dụng ensemble/hàm mất mát chuyên biệt cho lớp thiểu số.

\subsection{Giao thức thực nghiệm mở rộng}

Outer train/test split tiếp tục được giữ cố định ở tỷ lệ 80/20 với \texttt{random\_state=42}. Test set giữ nguyên phân bố ban đầu, không được resample và không tham gia lựa chọn mô hình hay threshold. Từ outer-training set, nhóm tách inner validation với seed 2026 để xác định operating threshold; một điểm vận hành chính yêu cầu Specificity trên validation đạt tối thiểu 85\%. Threshold sau đó được cố định trước khi đánh giá trên outer test.

Các chỉ số báo cáo gồm Accuracy, Precision, Recall, F1, MCC và ROC-AUC. Trong bối cảnh mất cân bằng, F1 và MCC được dùng để đánh giá sự cân bằng tổng thể, còn Recall được xem là guardrail cho khả năng phát hiện stroke. Việc chọn threshold trên validation thay vì hậu nghiệm trên test nhằm hạn chế thiên lệch do lựa chọn điểm vận hành sau khi đã quan sát kết quả cuối.

\subsection{Dataset 1: so sánh các họ mô hình mới}

CatBoost được huấn luyện trực tiếp trên các biến phân loại nguyên bản với \texttt{auto\_class\_weights=Balanced}, không One-Hot Encoding, PCA hay SMOTE. LightGBM giữ toàn bộ training set và dùng \texttt{scale\_pos\_weight}. Balanced Random Forest tạo các mẫu cân bằng bên trong quá trình xây dựng cây. EasyEnsemble và Focal-loss XGBoost được dùng như các thí nghiệm sàng lọc bổ sung cho bài toán mất cân bằng.

\begin{table}[H]
\centering
\scriptsize
\caption{Dataset 1: kết quả các phương pháp mở rộng trên cùng outer test}
\label{tab:d1_new_methods}
\resizebox{\linewidth}{!}{%
\begin{tabular}{lrrrrrrp{3.6cm}}
\toprule
\textbf{Phương pháp} & \textbf{Acc.} & \textbf{Prec.} & \textbf{Recall} & \textbf{F1} & \textbf{MCC} & \textbf{ROC-AUC} & \textbf{Vai trò}\\
\midrule
XGBoost cải tiến (Exp. 13) & 0,8434 & 0,1944 & 0,7000 & 0,3043 & 0,3119 & 0,8241 & Mốc cải tiến trước đó\\
CatBoost & 0,7867 & 0,1640 & \textbf{0,8200} & 0,2733 & 0,3036 & \textbf{0,8364} & Recall/ROC-AUC cao\\
LightGBM weighted & 0,8327 & 0,1543 & 0,5400 & 0,2400 & 0,2220 & 0,7865 & Không cải thiện\\
\textbf{Balanced Random Forest} & \textbf{0,8483} & \textbf{0,2000} & 0,7000 & \textbf{0,3111} & \textbf{0,3184} & 0,8274 & \textbf{Tốt nhất tổng thể}\\
EasyEnsemble$^{*}$ & 0,8454 & 0,1860 & 0,6400 & 0,2883 & 0,2860 & 0,8309 & Screening\\
Focal-loss XGBoost & 0,8405 & 0,1844 & 0,6600 & 0,2882 & 0,2893 & 0,8317 & Không cải thiện\\
\bottomrule
\end{tabular}}
\end{table}

Balanced Random Forest tạo ra điểm vận hành cân bằng nhất. So với XGBoost cải tiến ở Experiment 13, Recall giữ nguyên 0,70, trong khi Accuracy tăng từ 0,8434 lên 0,8483, Precision từ 0,1944 lên 0,2000, F1 từ 0,3043 lên 0,3111 và MCC từ 0,3119 lên 0,3184. Mức tăng tương đối khoảng 2,2\% ở F1 và 2,1\% ở MCC, vì vậy được diễn giải là \textbf{cải thiện biên (marginal improvement)}, không phải sự vượt trội lớn.

CatBoost cho Recall và ROC-AUC cao nhất, lần lượt 0,82 và 0,8364. Điều này cho thấy khả năng xếp hạng nguy cơ tốt khi xử lý trực tiếp biến categorical, nhưng lợi thế ranking chưa chuyển thành F1/MCC tốt nhất tại operating point được chọn. LightGBM, EasyEnsemble và Focal-loss XGBoost không vượt Balanced Random Forest hoặc XGBoost cải tiến trên các chỉ số cân bằng chính.

$^{*}$EasyEnsemble sử dụng cấu hình giới hạn tài nguyên tính toán sau khi cấu hình mặc định không phù hợp với môi trường chạy; vì vậy kết quả được xem là screening experiment, không phải bằng chứng chính để xếp hạng mô hình.

\subsection{Dataset 2: class weighting so với data-level balancing}

Dataset 2 sau làm sạch có 5.542.261 quan sát; outer-training set gồm 4.433.808 và outer-test set gồm 1.108.453 mẫu. LightGBM được chạy trên toàn bộ outer-training set, không undersampling hay SMOTE; mất cân bằng được xử lý bằng \texttt{scale\_pos\_weight}. Kết quả được so sánh với cấu hình XGBoost + undersampling + PCA tốt nhất ở Experiment 13.

\begin{table}[H]
\centering
\scriptsize
\caption{Dataset 2: so sánh XGBoost cải tiến và LightGBM weighted trên full data}
\label{tab:d2_new_methods}
\resizebox{\linewidth}{!}{%
\begin{tabular}{lrrrrrr}
\toprule
\textbf{Phương pháp} & \textbf{Acc.} & \textbf{Prec.} & \textbf{Recall} & \textbf{F1} & \textbf{MCC} & \textbf{ROC-AUC}\\
\midrule
\textbf{XGBoost + undersampling + PCA} & \textbf{0,9847} & \textbf{0,7401} & 0,9886 & \textbf{0,8465} & \textbf{0,8484} & \textbf{0,9987}\\
LightGBM weighted & 0,9587 & 0,5079 & \textbf{0,9926} & 0,6720 & 0,6943 & 0,9973\\
\bottomrule
\end{tabular}}
\end{table}

LightGBM tăng Recall nhẹ từ 0,9886 lên 0,9926 nhưng Precision giảm mạnh từ 0,7401 xuống 0,5079; F1 giảm còn 0,6720 và MCC còn 0,6943. Kết quả cho thấy class weighting giúp hạn chế false negative nhưng tạo nhiều false positive hơn. Trên Dataset 2, \textbf{data-level balancing bằng undersampling kết hợp XGBoost cho trade-off Precision--Recall tốt hơn algorithm-level class weighting của LightGBM}.

CatBoost cũng được khảo sát cho Dataset 2, nhưng huấn luyện full 4,43 triệu quan sát vượt giới hạn bộ nhớ của môi trường thực nghiệm. Một probe dùng stratified subset được thực hiện để kiểm tra tính khả thi, song không được đưa vào bảng so sánh full-data và không được dùng để kết luận hơn/kém so với XGBoost hay LightGBM. Cách trình bày này tránh đánh đồng một thí nghiệm giới hạn tài nguyên với các experiment sử dụng toàn bộ training set.

\subsection{Đúc kết từ thử nghiệm mở rộng}

Các thử nghiệm cho thấy thay thế XGBoost bằng một họ thuật toán mới không tự động tạo ra cải thiện. Với Dataset 1 nhỏ và mất cân bằng mạnh, Balanced Random Forest cho cải thiện nhỏ nhưng nhất quán về Accuracy, Precision, F1 và MCC khi giữ nguyên Recall. CatBoost lại nổi bật ở Recall và ROC-AUC, cho thấy một operating objective khác có thể dẫn đến lựa chọn mô hình khác.

Với Dataset 2 quy mô lớn, XGBoost kết hợp undersampling và PCA vẫn là cấu hình tốt nhất xét đồng thời Precision, Recall, F1 và MCC. LightGBM weighted đạt Recall cao hơn nhưng phải đánh đổi bằng số false positive lớn hơn. Do đó, ảnh hưởng của \textbf{cơ chế xử lý mất cân bằng và operating threshold ít nhất quan trọng tương đương việc lựa chọn họ mô hình}.

Kết luận của giai đoạn mở rộng không phải là một thuật toán mới có thể thay thế XGBoost trong mọi trường hợp. Thay vào đó, phương pháp tối ưu phụ thuộc vào cấu trúc và quy mô dữ liệu: \textbf{Balanced Random Forest là phương án mở rộng tốt nhất cho Dataset 1}, trong khi \textbf{XGBoost + undersampling + PCA tiếp tục là cấu hình tốt nhất cho Dataset 2}. Các kết quả âm tính của LightGBM, EasyEnsemble và Focal-loss XGBoost cũng được giữ lại để bảo đảm báo cáo phản ánh đầy đủ quá trình thử nghiệm thay vì chỉ lựa chọn hậu nghiệm các kết quả thuận lợi.

% Trigger CI rebuild after adding the PDF workflow.
